\PassOptionsToPackage{table}{xcolor}
\documentclass[aspectratio=169]{beamer}

\usetheme{Madrid}
\usecolortheme{whale}

% Remove navigation symbols
\setbeamertemplate{navigation symbols}{}
\setbeamertemplate{blocks}[rounded][shadow=false]

% --- Custom Design ---
\definecolor{polytechblue}{RGB}{30, 60, 120}
\setbeamercolor{structure}{fg=polytechblue}
\setbeamercolor{title}{fg=white, bg=polytechblue}
\setbeamercolor{block title}{bg=polytechblue, fg=white}
\setbeamerfont{title}{series=\bfseries}

\usepackage[utf8]{inputenc}
\usepackage[T1]{fontenc}
\usepackage[french]{babel}
\usepackage{graphicx}
\usepackage{booktabs}
\usepackage{tikz}
\usepackage{multicol}
\usepackage{colortbl}

% Title Slide Info
\title{Multimedia Image Retrieval}
\subtitle{Recherche d’Images Unimodale (Cars) et Multimodale (CLIP)}
\author[Alla, Xu, Toussaint]{Julien XU \and Abdelouahad Alla \and Paul TOUSSAINT}
\institute[UMONS]{University of Mons \\ Faculté Polytechnique \\ Faculté des Sciences}
\date{Juin 2026}

\begin{document}

% 1. Title Page
\begin{frame}[plain]
    \titlepage
\end{frame}

% 2. Table of Contents
\begin{frame}{Sommaire}
    \tableofcontents
\end{frame}

% --- Section 1: Introduction ---
\section{Introduction}

\begin{frame}{Introduction et Objectifs}
    \begin{block}{Contexte du projet}
        Développement d'un système de recherche d'information multimédia appliqué au domaine automobile (Cars) et aux scènes quotidiennes (Flickr8K).
    \end{block}
    
    \begin{block}{Trois volets majeurs}
        \begin{itemize}
            \item \textbf{Unimodal} : Recherche image-vers-image sur 10 000 voitures.
            \item \textbf{Multimodal} : Recherche croisée texte/image via CLIP.
            \item \textbf{Déploiement} : Solution Cloud optimisée sur Render.
        \end{itemize}
    \end{block}
\end{frame}

% --- Section 2: Unimodale ---
\section{Recherche Unimodale (Cars)}

\begin{frame}{Dataset et Architecture Unimodale}
    \begin{columns}[T]
        \begin{column}{0.6\textwidth}
            \begin{itemize}
                \item \textbf{Données} : Base "Cars" de 10 000 images (5 classes cibles).
                \item \textbf{Descripteurs exploités} :
                \begin{itemize}
                    \item \textbf{Classiques} : BGR, HSV, HOG, LBP, SIFT/ORB pooled.
                    \item \textbf{Deep} : ResNet50 et ViT-B/16.
                \end{itemize}
                \item \textbf{Distance en mode "Auto"} :
                \begin{itemize}
                    \item $\chi^2$ pour les histogrammes.
                    \item Euclidienne pour le reste.
                    \item Cosinus pour le Deep Learning.
                \end{itemize}
            \end{itemize}
        \end{column}
        \begin{column}{0.4\textwidth}
            \centering
            \includegraphics[width=0.8\textwidth]{figures/interface_unimodale.png}\\
            \tiny Interface Flask (Fig 2.1)
        \end{column}
    \end{columns}
\end{frame}

\begin{frame}{Fusion de Descripteurs : L'approche hybride}
    \begin{block}{L'intérêt de la fusion}
        Combiner plusieurs axes (ex: Couleur + Deep Learning) pour améliorer la précision.
    \end{block}
    \begin{itemize}
        \item \textbf{Fusion tardive} : On calcule les distances pour chaque descripteur séparément.
        \item \textbf{Normalisation Min-Max} : Ramener les distances entre 0 et 1.
        \item \textbf{Score final} : Moyenne des distances normalisées.
    \end{itemize}
    \begin{exampleblock}{Exemple concret}
        Fusion \textbf{HSV} (Couleur) + \textbf{ViT} (Sémantique) pour compenser les erreurs de forme.
    \end{exampleblock}
\end{frame}

\begin{frame}{Résultats Unimodaux : Quantitative}
    \begin{table}[]
        \centering
        \resizebox{0.8\textwidth}{!}{
        \begin{tabular}{lcccc}
            \toprule
            \textbf{Descripteur} & \textbf{P@50} & \textbf{mAP@50} & \textbf{R-Prec} & \textbf{vs HSV}\\ \midrule
            \rowcolor{blue!15} \textbf{ViT-B/16} & \textbf{0.231} & \textbf{0.096} & \textbf{0.151} & \textbf{+134\%} \\
            ResNet50 & 0.172 & 0.070 & 0.128 & +70\% \\
            HSV (Baseline) & 0.143 & 0.041 & 0.105 & - \\
            \bottomrule
        \end{tabular}
        }
    \end{table}
    \begin{block}{Impact du Deep Learning}
        Le passage aux Transformers multiplie par \textbf{plus de 2} la précision moyenne par rapport aux méthodes classiques couleur.
    \end{block}
\end{frame}

\begin{frame}{Efficacité Technique : Compromis Temps / Mémoire}
    \begin{table}[]
        \centering
        \resizebox{0.9\textwidth}{!}{
        \begin{tabular}{lccc}
            \toprule
            \textbf{Descripteur} & \textbf{Dim. Vecteur} & \textbf{Indexation (10k)} & \textbf{Recherche / req} \\ \midrule
            HSV (Rapide) & 540 & \textbf{26.7 s} & 3.70 $\mu$s \\
            ResNet50 & 2048 & 594.5 s & 5.86 $\mu$s \\
            \rowcolor{blue!5} ViT-B/16 & 768 & 1364.5 s & 2.21 $\mu$s \\
            HOG (Lourd) & \textbf{8100} & 27.9 s & 7.00 $\mu$s \\
            \bottomrule
        \end{tabular}
        }
    \end{table}
    \begin{itemize}
        \item \textbf{ViT} : Meilleur équilibre entre taille de vecteur et précision.
        \item \textbf{HOG} : Très lent à la recherche commerciale à cause de sa dimension.
        \item \textbf{HSV} : Idéal pour une indexation ultra-légère.
    \end{itemize}
\end{frame}

\begin{frame}{Analyse Qualitative : Cas de succès et limites}
    \begin{columns}[T]
        \begin{column}{0.5\textwidth}
            \centering
            \includegraphics[width=0.9\textwidth]{figures/resultats_unimodal_succes.png}\\
            \tiny Succès : Peugeot Rifter (Fig 2.2)
        \end{column}
        \begin{column}{0.5\textwidth}
            \centering
            \includegraphics[width=0.9\textwidth]{figures/resultats_unimodal_echec.png}\\
            \tiny Limites : Peugeot 508 (Fig 2.5)
        \end{column}
    \end{columns}
    \vspace{0.3cm}
    \begin{alertblock}{Le "Semantic Gap"}
        Similarité visuelle (forme/couleur) $\neq$ Appartenance à la marque.
    \end{alertblock}
\end{frame}

% --- Section 3: Multimodale ---
\section{Recherche Multimodale (CLIP)}

\begin{frame}{Modèle CLIP et Dataset Flickr8K}
    \begin{itemize}
        \item \textbf{Dataset} : 8 000 images quotidiennes, 40 000 captions (5/image).
        \item \textbf{Modèle} : \textbf{CLIP ViT-B/32} (Contrastive Language-Image Pretraining).
        \item \textbf{Architecture} : Espace latent commun Image/Texte de dimension 512.
    \end{itemize}
    \begin{block}{Recherche ultra-rapide via FAISS}
        Indexation des 40 000 captions dans un index \textbf{IndexFlatIP} pour un temps de réponse instantané.
    \end{block}
\end{frame}

\begin{frame}{Évaluation Multimodale : Résultats CLIP}
    \begin{columns}
        \begin{column}{0.5\textwidth}
            \centering
            \includegraphics[width=0.9\textwidth]{figures/courbe_text_to_image.png}\\
            \tiny P-R Text-to-Image (Fig 3.2)
        \end{column}
        \begin{column}{0.5\textwidth}
            \centering
            \includegraphics[width=0.9\textwidth]{figures/courbe_image_to_text.png}\\
            \tiny P-R Image-to-Text (Fig 3.3)
        \end{column}
    \end{columns}
    \begin{block}{Performance}
        mAP Text-to-Image de \textbf{0.4969}. Le modèle capture extrêmement bien les descriptions en langage naturel (\textit{Zero-shot}).
    \end{block}
\end{frame}

% --- Section 4: Cloud ---
\section{Hébergement et Ingénierie}

\begin{frame}{Déploiement Cloud sur Render}
    \begin{itemize}
        \item \textbf{Stack technique} : Docker + Flask API + Render (Platform as a Service).
        \item \textbf{Contrainte RAM (512 Mo)} :
        \begin{itemize}
            \item Impossible de charger le modèle CLIP complet sans plantage (Erreur 502).
        \end{itemize}
    \end{itemize}
    \begin{block}{Solutions d'optimisation}
        \begin{itemize}
            \item \textbf{Empreinte mémoire} : Chargement différé et exploitation de modèles allégés.
            \item \textbf{Offline Processing} : Envoi des embeddings pré-calculés (.npy) sur le serveur.
        \end{itemize}
    \end{block}
\end{frame}

% --- Conclusion ---
\section{Conclusion}

\begin{frame}{Conclusion et Bilan}
    \begin{block}{Points forts du travail}
        \begin{itemize}
            \item Système \textbf{hybride} performant (Classique + Deep).
            \item Maîtrise des modèles \textbf{state-of-the-art} (CLIP, ViT).
            \item Solution logicielle complète (Backend, Frontend, Cloud).
        \end{itemize}
    \end{block}
    
    \begin{block}{Perspectives}
        \begin{itemize}
            \item Fine-tuning des modèles sur des données spécifiques (Automobile).
            \item Extension à des bases de données de plus grande échelle (Mil-scale).
        \end{itemize}
    \end{block}
    \centering \vfill \textbf{Merci pour votre attention ! Des questions ?}
\end{frame}

\end{document}
